\chapter{第五章}
\section{紧集与进空间的等价命题}
\begin{theorem}[在拓扑空间 $X$ 中，子集 $A$ 紧 $\Longleftrightarrow$ $A$ 的每个开覆盖都有有限子覆盖。]
$X$ 是拓扑空间，$A\subset X$，$A$ 是紧集当且仅当对 $X$ 的任意一开集族 $\mathcal{U}$，若 $A\subset\bigcup\mathcal{U}$，则存在有限子族 $\mathcal{U}_1\subset\mathcal{U}$，使得 $A\subset\bigcup\mathcal{U}_1$。
\end{theorem}
\begin{example}闭区间 $[0,1]$ 是紧的
在拓扑空间 $\mathbb{R}$（带通常拓扑）里取子集
\[
A=[0,1].
\]

随便给它一个“乱七八糟”的开覆盖 $\mathcal U$，例如
\[
\mathcal U=
\left\{
\left(-\frac12,\frac34\right),
\left(\frac12,\frac32\right),
\left(-1,2\right),
\left(0,\frac13\right),
\left(\frac23,1.5\right),
\ldots
\right\}.
\]
显然
\[
[0,1]\subset \bigcup_{U\in\mathcal U} U,
\]
因为这些开区间的并至少把整个 $[0,1]$ 都盖住了。

根据“紧 ⇔ 每个开覆盖有有限子覆盖”的定理，一定能从这堆开集里挑出有限个开集就覆盖 $[0,1]$。例如取
\[
U_1=\left(-\tfrac12,\tfrac34\right),\qquad
U_2=\left(\tfrac12,\tfrac32\right),
\]
就已经有
\[
[0,1]\subset U_1\cup U_2.
\]

这就是一个“开覆盖 $\mathcal U$ 有有限子覆盖 $\{U_1,U_2\}$”的具体例子，说明闭区间 $[0,1]$ 满足定理中的条件，因此是紧集。
\end{example}
\begin{theorem}[紧+连续映射，像集也是紧的]
若 $X$ 是紧拓扑空间，$f:X\to Y$ 是连续映射，则 $f(X)$ 是 $Y$ 的紧子集。
\end{theorem}
\begin{theorem}[紧空间判定-闭集族有限交不是空的]
设 $X$ 是拓扑空间，则 $X$ 为紧空间的充要条件是：  
对 $X$ 中每个具有有限交性质的闭集族 $\mathscr{F}$，都有
\[
\bigcap \mathscr{F} \ne \varnothing.
\]
\end{theorem}
\begin{theorem}[闭子空间继承紧]
若 $X$ 是紧空间，则 $X$ 的每个闭子空间也是紧空间。
\end{theorem}
\section{R和$R^n$中紧集}
\begin{itemize}
  \item 在研究 $\mathbb{R}$ 中的紧集时，先给出“有界”的概念：  
  若 $A\subset\mathbb{R}$，且存在 $a\in\mathbb{R}$、$a>0$，使得
  \[
    A\subset[-a,a],
  \]
  则称集合 $A$ 是\textbf{有界集}。
  

  \item 结论 5.5：令 $X=\mathbb{R}$，取通常拓扑。若 $A\subset\mathbb{R}$ 是紧集，则 $A$ 是有界集。

  \item 结论 5.6：令 $X=\mathbb{R}$，取通常拓扑。若 $[a,b]$ 是 $\mathbb{R}$ 中的闭区间，则 $[a,b]$ 是紧集。
\end{itemize}
\begin{theorem}[R中紧=有界闭集]
设 $X=\mathbb{R}$，在 $X$ 上取通常拓扑，$A\subset X$。则 $A$ 是紧集的充要条件是：$A$ 是 $\mathbb{R}$ 中的有界闭集。
\end{theorem}

\begin{theorem}[紧的空间乘积也是紧空间]
若 $X$ 与 $Y$ 都是紧空间，则积空间 $X\times Y$ 也是紧空间。
\end{theorem}
\begin{corollary}[在通常拓扑下，欧氏空间 $\mathbb{R}^n$ 里的有界闭盒子 $[-a,a]^n$ 总是紧集。]
对任意 $n\in\mathbb{N}$ 及 $a\in\mathbb{R}$ 且 $a>0$，集合
\[
[-a,a]^n \subset \mathbb{R}^n
\]
是 $\mathbb{R}^n$ 中的紧集。
\end{corollary}
\begin{definition}[Hausdorff 空间 / $T_2$ 空间]
设 $(X,\mathcal{T})$ 为拓扑空间。若对任意两个不同的点 $x,y\in X$，都存在不相交的开集
$U,V\in\mathcal{T}$，使得
\[
x\in U,\qquad y\in V,\qquad U\cap V=\varnothing,
\]
则称 $X$ 为 $T_2$ 拓扑空间，简称为 $T_2$ 空间，又称 Hausdorff 空间。

因此，带通常拓扑的欧氏空间 $\mathbb{R}^n$ 是一个 $T_2$ 拓扑空间。
\end{definition}
对于 $x,y\in\mathbb{R}^n$，记
\[
x=(x_i:i\le n),\qquad y=(y_i:i\le n),
\]
定义
\[
d(x,y)=\sqrt{\sum_{i=1}^{n}(x_i-y_i)^2}.
\]
则 $d(x,y)$ 具有如下性质：
\begin{itemize}
  \item[(1)] $d(x,y)=0$ 当且仅当 $x=y$；
  \item[(2)] $d(x,y)=d(y,x)$；
  \item[(3)] $d(x,y)\le d(x,z)+d(z,y),\qquad x,y,z\in\mathbb{R}^n$。
\end{itemize}
\begin{theorem}[T2空间紧集必是闭集]
如果 $X$ 是 $T_2$ 拓扑空间，$A\subset X$ 是 $X$ 中的紧集，则 $A$ 是 $X$ 中的闭集。
\end{theorem}
\begin{theorem}[$R^n$中紧和有界闭等价]
设 $A\subset \mathbb{R}^n$，则 $A$ 为紧集的充要条件是：$A$ 是 $\mathbb{R}^n$ 中的有界闭集。
\end{theorem}
\begin{theorem}[紧到T2+连续映射=闭]
如果 $X$ 是紧空间，$f:X\to Y$ 是连续映射，$Y$ 是 $T_2$ 空间，则 $f(X)$ 是 $Y$ 中的闭集。
\end{theorem}
\begin{theorem}[连续映射，X到R，X紧子集上有映射极值的原像]
如果 $f:X\to \mathbb{R}$ 是连续映射，且 $A\subset X$ 是 $X$ 中的紧集，则 $f$ 在 $A$ 上可取得最大（小）值。
\end{theorem}
\begin{proof}
因为 $A$ 是 $X$ 中的紧集，而 $f:X\to\mathbb{R}$ 是连续映射，则 $f|_{A}:A\to\mathbb{R}$ 
也是连续映射。

由连续函数在紧空间上取得最大值和最小值的基本定理可知，存在点
\[
x_{\max},\,x_{\min}\in A,
\]
使得
\[
f(x_{\max})=\max_{x\in A} f(x),\qquad 
f(x_{\min})=\min_{x\in A} f(x).
\]

因此，$f$ 在紧集 $A$ 上能够取得其最大值和最小值。
\end{proof}
\section{紧空间的无限积空间}
\begin{lemma}[任何具有有限交性质的集族都可以扩张成一个包含它的、具有有限交性质的极大集族]
设 $X$ 是拓扑空间，$\mathcal{F}$ 是 $X$ 中的具有有限交性质的集族，
则存在 $\mathcal{U}_{\mathcal{F}}$ 是具有有限交性质的集族，使得 
$\mathcal{F} \subset \mathcal{U}_{\mathcal{F}}$，且对任一具有有限交性质的集族 
$\mathcal{V}_{\mathcal{F}}$，若 $\mathcal{U}_{\mathcal{F}} \subset \mathcal{V}_{\mathcal{F}}$，
则 $\mathcal{U}_{\mathcal{F}} = \mathcal{V}_{\mathcal{F}}$
（称 $\mathcal{U}_{\mathcal{F}}$ 是包含 $\mathcal{F}$ 且具有有限交性质的极大族）。
\end{lemma}
\begin{example}[在 $\mathbb{R}$ 上理解引理 5.1]

考虑集族
\[
\mathcal{F}=\{[0,1],\ [0,\tfrac12],\ [0,\tfrac13],\dots\}.
\]

\begin{itemize}
    \item[\color{purple}{(1)}] {\color{purple}{有限交性质}}：  
    任意有限多个 $[0,1/n]$ 的交仍为非空，例如
    \[
    [0,1/2]\cap[0,1/5]\cap[0,1/7]=[0,1/7]\neq\varnothing.
    \]

    \item[\color{purple}{(2)}] {\color{purple}{构造极大有限交族}}：  
    一个包含 $\mathcal{F}$ 的集族为
    \[
    \mathcal{U_F}=\{A\subseteq\mathbb{R}\mid 0\in A\},
    \]
    因为所有 $[0,1/n]$ 都含 $0$。

    \item[\color{purple}{(3)}] {\color{purple}{为什么它是极大？举具体例子说明不能再加入集合}}：
        \begin{itemize}
            \item 所有 $A\in\mathcal{U_F}$ 都含 $0$，因此任意有限交均含 $0$，故非空。
            \item 试图加入集合 $B=(1,2)$（显然 $0\notin B$）。  
            因为 $[0,1]\in\mathcal{U_F}$，所以
            \[
            B\cap[0,1]=\varnothing,
            \]
            有限交性质被破坏，因此 $B$ 不能加入。这里[0,1]是特意选出来的。
            \item 所以任何不含 $0$ 的集合都无法加入；  
            而含 $0$ 的集合已经全部在 $\mathcal{U_F}$ 里，故它是极大的。
        \end{itemize}
\end{itemize}

\end{example}
\begin{theorem}[任意多个紧空间的乘积仍然是紧的]
对 \(\alpha \in \Lambda\)，若 \(X_\alpha\) 是紧空间，则 \(\prod_{\alpha \in \Lambda} X_\alpha\) 也是紧空间。
\end{theorem}
\section{完备映射}
\begin{definition}[完备映射=连续闭满+像中所有点的原像集都是紧的]
设 $f:X\to Y$ 是连续的闭满映射，若对任意 $y\in Y$，$f^{-1}(y)$ 是 $X$ 中的紧子集，则称 $f$ 是完备映射。
\end{definition}
\begin{theorem}[完备映射,像是紧的原像才是紧的]
若 $f:X\to Y$ 是完备映射，则 $X$ 是紧空间当且仅当 $Y$ 是紧空间。
\end{theorem}
\begin{theorem}[完备映射保持林德洛夫空间]
若 $f:X\to Y$ 是完备映射，则 $X$ 是 Lindelöf 空间当且仅当 $Y$ 是 Lindelöf 空间。
\end{theorem}

\begin{theorem}[拓扑*紧映射到自己是完备的]
若 $X$ 是拓扑空间、$Y$ 是紧空间，则投影映射 
\[
P_X : X\times Y \to X
\]
是完备映射。
\end{theorem}
\begin{corollary}[林德洛夫*紧=林德洛夫]
若 $X$ 是 Lindelöf 空间，$Y$ 是紧空间，则 $X\times Y$ 是 Lindelöf 空间。
\end{corollary}

\begin{theorem}[完备保持第二可数]
若 $f:X\to Y$ 是完备映射，$X$ 是第二可数空间，则 $Y$ 也是第二可数空间。
\end{theorem}

\section{第一纲集与第二纲集}
\begin{definition}[第一纲集与第二纲集]
设 $X$ 是拓扑空间，$A\subset X$。如果 $\overline{A}^{\,\circ}=\varnothing$，则称 $A$ 是 $X$ 中的无处稠密集。

如果一集合 $A$ 是可数个无处稠密集的并，则称 $A$ 是第一纲集；如果 $A$ 不是第一纲集，则称 $A$ 是第二纲集。
\end{definition}
\begin{note}
     nowhere dense set 疏集就是无处稠密集英文是一个
\end{note}
\begin{definition}[疏集与第一/第二纲集]
设 $E$ 是度量空间 $X$ 中的点集。  
如果 $E$ 不在 $X$ 的任何开集中稠密，则称 $E$ 是\textbf{疏集}。  

如果一个集合可以表示为可数个疏集的并集，称为\textbf{第一纲 (meager set) 集}；  
否则称为\textbf{第二纲集}。
\end{definition}
集合 $E$ 为疏集，当且仅当 $\overline{E}$ 不能充满 $X$ 中的任何一个开集，
即
$
\operatorname{int}(\overline{E}) = \varnothing.
$int是 interior（内点集/内部） 的简写。

意思是：集合 A 的内部，也就是包含在 A 中的最大开集。
\begin{example}
$\mathbb{R}$ 中的有理数集 $\mathbb{Q}$ 是第一纲集，无理数集 $I$ 是第二纲集。
\end{example}